\section{Methods: four RL algorithms}

We solve the same MDP with four algorithms, chosen so that each pair of them contrasts one idea from the course~\citep{sutton2018}.

\begin{table}[h]
\centering
\caption{The four algorithms and their final settings.}
\label{tab:algos}
\small
\begin{tabularx}{\linewidth}{@{}l>{\raggedright}p{2.4cm}>{\raggedright}p{2.8cm}X@{}}
\toprule
Algorithm & Course topic & State & Key details \\
\midrule
Value Iteration & Dynamic programming, Bellman equations & Tabular: $s$ (1--8) $\times$ 10 bins of $q$ = 80 states & Transition probabilities and expected rewards counted from the logs (every logged position as a round start), then Bellman optimality updates until the residual is below $10^{-9}$ (about 8 sweeps) \\
Q-learning & Temporal-difference learning & Same 80 tabular states & $Q(s,a) \leftarrow Q(s,a) + \alpha[r + \max_{a'} Q(s',a') - Q(s,a)]$, $\alpha=0.05$, $\varepsilon$-greedy ($1\to0.05$ over 80\% of training), 10\,000 episodes \\
DQN~\citep{mnih2015} & Function approximation, deep RL & Continuous features $[s/8, p, H/5, q, a]$ & MLP 5$\to$64$\to$64$\to$2, replay buffer of 50k steps, batch 128, update every 4 steps, target network synced every 500 steps, Huber loss, Adam ($10^{-3}$), 3000 episodes \\
REINFORCE~\citep{williams1992} & Policy gradient & Same continuous features & Softmax policy MLP 5$\to$64$\to$2, Monte Carlo returns, mean-return baseline, advantages divided by their standard deviation, Adam ($10^{-3}$), 3000 episodes \\
\bottomrule
\end{tabularx}
\end{table}

\paragraph{Episodes.} All four algorithms treat each draft$\to$verify round as one short episode of at most
8 decisions with $\gamma=1$: \CONT{} bootstraps from the next state of the same round, and \STOP{} ends it.
An answer has a fixed number of tokens whatever the policy, so a whole-answer return would only measure
time and $\lambda$ would have no effect. Per round, maximising
$\mathbb{E}[\text{gained}-\lambda\,\text{cost}]$ is exactly Dinkelbach's form of maximising throughput.
Steps where \STOP{} is forced ($s=8$, end of the answer) are stored as \STOP{} and get no policy gradient,
and \CONT{} at $s=8$ is excluded from bootstrap targets.

\paragraph{Choosing $\lambda$.} Each algorithm is trained twice: first with $\lambda$ equal to the best
fixed-$k$ throughput on the training prompts, then with $\lambda$ equal to the throughput its own first policy
reached (one Dinkelbach update).

\paragraph{Comparison axes.}
\begin{itemize}
  \item Model-based (Value Iteration) vs model-free (Q-learning).
  \item Tabular (Q-learning) vs function approximation (DQN).
  \item Value-based (DQN) vs policy-based (REINFORCE).
\end{itemize}

\paragraph{Baselines (not RL).}
\begin{itemize}
  \item Target model alone, with no speculation (speedup $= 1.0\times$ by definition).
  \item Fixed draft length $k = 1$ to $8$.
  \item Confidence threshold: stop when $p < 0.4$, as in HuggingFace's dynamic speculation~\citep{hfdynamic}.
\end{itemize}
